% Ad Atticum 5.20 --- headnote
% ad-atticum-05-20, between 21 and 29 December 51 BC, Cilicia (probably Tarsus)

\textit{Cicero to Atticus, written somewhere in Cilicia
--- probably at Tarsus, on the way back from Pindenissus
to winter quarters at Laodicea --- between the twelfth
and the fourth days before the Kalends of January
(21--29 December) 51 BC (the manuscript dateline:
\textit{Scr. in Cilicia inter a. d. xii et iv K. Ian.
a. 703 (51)}). The long closing letter of \textit{Ad
Atticum} book 5: the work of an entire proconsular year
gathered into one despatch and sent off to the man for
whom the whole correspondence is intended as its first
audience. The book that began at Brundisium in May
with the dread of having to govern at all ends here at
Pindenissus in December with the army being demobilised
into winter quarters and the writer composing the
official letter to the Senate that will follow.}

\textit{The letter is structurally a composite. The
opening section is the bare military headline: the fall
of Pindenissus on the morning of the Saturnalia after a
fifty-seven day investment, with the jocular aside that
its name will mean nothing to Atticus and the
self-mocking acknowledgement (\textit{en epitome}) that
this whole campaign cannot be made to sound like
Aetolia. Sections 2 to 4 then back-fill the narrative
from Iconium through the Amanus to Issus to Pindenissus,
turning aside in section 4 to settle the rival Bibulus
in a single famous one-liner (\textit{coepit loreolam
in mustaceo quaerere} --- ``he began to look for his
little laurel in a baker's tart''). The Saturnalia
plunder is distributed to the soldiers --- the horses
kept back --- and the slaves auctioned off on the third
day of the feast. Section 6 is the great private
passage: Cicero on his own integrity in the province,
not as continentia but as the deepest personal pleasure
he has ever known --- ``I did not know myself, and did
not know well enough what I could do in this line.''
The Greek-tag scaffolding (\textit{Momus}, \textit{en
epitome}, \textit{pephusiomai}, \textit{lampra},
\textit{aprositon}, \textit{adorodoketon}) is dense,
intimate, and confidential: the inside-Atticus voice
Cicero uses with no one else. Sections 7 to 10 swing
out to Roman news --- Caesar's prospects, the new
political shape Atticus has just reported, Quintus the
nephew's coming-of-age toga, the wardship of young
Brutus, the household at Ephesus, and the standing
business of the boy in the house of Pammenes. The
private letter does what the public one will not be
able to: ledger the year as a whole.}
